\subsection{최종 답변의 자동 품질과 자원 사용 특성 (RQ3)}

\label{sec:ch4_4}

Jev의 Token F1은 46.20점에서 47.31점으로, LF-Eval은 67.57점에서 70.15점으로 높아졌다. 차이의 CI는 각각 [+0.02, +2.27]점과 [+0.37, +4.79]점이었다. Token F1은 정답과 생성 답변의 공백 단위 표현 중복을 측정하며, LF-Eval은 정답 답변과 gold에 비추어 법률적 결론·누락·모순을 평가한다. 두 평가에는 공식 코드의 앞 800자 처리 규칙을 적용하였다. Token F1의 별도 centered-bootstrap 양측 p값은 약 0.0525로, percentile CI와 검정의 구성 방식이 다르므로 모든 자동 지표가 유의하게 개선되었다고 해석하지 않았다.

선택 문맥은 평균 203.67토큰에서 215.60토큰으로, 전체 API 입력은 2302.67토큰에서 2314.58토큰으로 늘었다. 출력은 72.16토큰과 73.98토큰이었다. 문맥 길이의 중앙값은 164토큰과 170.5토큰이며, 문항별 길이 차이의 중앙값은 0이었다. 선택 목록은 양쪽 모두 478개로 동일하므로, 이는 선택 수 감소가 아닌 선택 내용 변화에 따른 결과이다. 본 비교에서 입력량 절감은 관찰되지 않았다. 사용량 토큰은 Token F1의 공백 단위와 다르며, 재사용 응답의 기록도 포함하므로 신규 청구액을 뜻하지 않는다.

추가로 저장된 선택 근거에 문턱값을 적용하는 오프라인 필터링을 수행하였다. 151개 의존 그룹의 5-fold 교차검증에서 학습 fold의 F1 하락이 1점 이내인 후보 중 근거를 가장 적게 남기는 문턱값을 선택하고 검증 fold에 적용하였다. scorer별 점수 척도를 고려해 동일한 크기의 분위수 격자를 사용하였다.

\begin{table}[htbp]
\centering
\caption{선택 후 필터링의 검증 fold 통합 결과(226문항).}
\small
\setlength{\tabcolsep}{3pt}
\begin{tabular}{@{}lrrrr@{}}
\hline
조건 & 유지 근거 수 & F1 & Precision & Recall \\
\hline
CE & 478 & 58.57 & 63.46 & 57.74 \\
CE+필터링 & 467 & 58.54 & 64.56 & 57.15 \\
Jev & 478 & 61.75 & 66.85 & 60.69 \\
Jev+필터링 & 426 & 61.43 & 69.80 & 58.26 \\
\hline
\end{tabular}
\end{table}

Jev는 근거의 89.12\%를 유지하면서 F1이 0.32점 낮아졌고, Precision은 2.95점 높아졌으나 Recall은 2.43점 낮아졌다(표 1). F1 차이 CI는 [$-$1.76, +1.07]점으로, 동등성이나 비열등성을 입증한 결과는 아니다. 근거가 하나도 남지 않은 문항은 CE 1개, Jev 4개였다. 같은 분포에서 근거를 90\% 유지한 기술통계에서도 F1은 CE 57.21, Jev 61.66이었다. 선택 전 필터링에서는 shortlist 후보를 약 절반으로 줄일 때 pool별 gold 유지 비율이 CE 66.14$\to$51.47\%, Jev 72.69$\to$69.30\%였다. 이는 pool별 집계로 그림 2의 문항별 평균 Recall과 분모가 다르다.

필터링 후 답변은 재생성하지 않았으므로 답변 품질 유지나 실제 비용 절감으로 해석하지 않았다. 또한 LF-Eval은 실제 제공 근거가 답변의 주장을 지지하는지를 직접 평가하지 않는다. 본 절의 결과는 자동 답변 품질과 자원 요구량에 한정된다.
